\section{Local State and the Journal}\label{sec:state}

\subsection{The data folder}\label{sec:folder}
All local state is in one folder, the per-user application-data folder of each system:
\begin{itemize}
\item Windows: \code{\%LOCALAPPDATA\%\textbackslash GhostVault}, which is local, not roaming, so the keys
do not travel with a roaming profile;
\item macOS: \path{~/Library/Application Support/GhostVault};
\item Linux: \path{$XDG_DATA_HOME/ghostvault}, or \path{~/.local/share/ghostvault}.
\end{itemize}
\lstref{lst:folder} shows its contents. On Linux and macOS, folders are created with mode 0700 and
files with 0600; on Windows the folder inherits the per-user permissions of the profile.

\begin{lstlisting}[caption={The data folder (modes on Linux and macOS).},label={lst:folder}]
<data folder>/              0700
  settings.json             0600  repositories, options
  manifests/<file id>.json  0600  one record per ghost
  thumbs/<file id>.gvp      0600  encrypted previews
  logs/ghostvault.log       0600  file ids, never names
  damaged/                  0700  unreadable files, kept
  vault.lock                0600  held while running
\end{lstlisting}

Nothing is staged on disk: chunks are encrypted and uploaded one at a time, and restored chunks go
straight into the hidden output file. Nothing in the folder records the folder's own location, so it
can be moved or restored from a backup. While GhostVault runs, it holds an exclusive lock on
\code{vault.lock} (\code{flock}, or \code{msvcrt.locking} on Windows); the operating system releases
it when the process ends, however it ends, and a second instance on the same folder is refused.

\subsection{Records}\label{sec:records}
Each ghost is one JSON record, \lstref{lst:record}, named after its random 128-bit file id. A record
is validated every time it is read: its shape, its version (5 or 6), and every value used to build a
request, choose a path or delete something must have the only form GhostVault creates: owner and
repository names, the base folder, chunk and preview paths, blob ids, nonces, the key length, the
cipher and associated-data label, the name of a temporary restore file and of the hidden name the
original is moved to, each beside its file. Records of format 6 must also have a chunk count that
matches their size. A record that fails is shown as damaged and left untouched; nothing acts on it,
and it affects no other ghost.

\begin{lstlisting}[language=gvjson,caption={A record of a ghost with ten chunks, abridged; the
values are from a real run against the simulated GitHub.},label={lst:record}]
{"format": "ghostvault-manifest", "version": 6,
 "fileId": "5ada355b920792c5ba14be70737ffc5a",
 "state": "ready", "displayName": "harbour.png",
 "original": {"name": "harbour.png", "size": 4860280,
   "path": "...", "sha256": "f9b10964...932b14",
   "mode": 420, "atimeNs": ..., "mtimeNs": ...,
   "identity": {"dev": ..., "ino": ..., "size": 4860280,
                "mtime_ns": ...}},
 "crypto": {"cipher": "AES-256-GCM",
            "key": "owrw+KXo...6508=",
            "aad": "ghostvault/v6"},
 "chunkSize": 524288, "padding": true,
 "chunks": [{"index": 0, "repo": "d825332f29363c8d",
   "path": "ghostvault/78/7e34107f...6b144794.bin",
   "plainSize": 524288, "nonce": "VG7sZuNj9KNv8fQb",
   "cipherSha256": "27b4368a...4e85e8",
   "blobSha": "b295c74a...616173be52d5"}, ...],
 "repos": {"d825332f29363c8d": {"owner": "alice",
   "repo": "vault1", "baseFolder": "ghostvault",
   "committed": true, "deleted": false}, ...},
 "preview": {"enabled": true, "remote": {"repo": ...,
   "path": "ghostvault/9c/962063f7...7c65b415.bin",
   "blobSha": ..., "committed": true}},
 "placement": {"strategy": "distributed",
               "eligible": 2, "repositories": 2},
 "notes": [], "restore": null}
\end{lstlisting}

\subsection{Durable writes}\label{sec:durable}
Every state file is written the same way: a new owner-only file is created beside the target
(\code{O\_EXCL}, \code{O\_NOFOLLOW}), written and flushed to disk, renamed over the old file, and
then the folder is flushed. After a crash the file holds either its old or its new content; a
leftover temporary file is removed at the next start. \tabref{tab:durable} lists what each step
means on each system. Flushing a file on macOS uses \code{F\_FULLFSYNC}, because \code{fsync} there
does not flush the drive's cache; Windows cannot flush a folder, so renames there are issued with
\code{MOVEFILE\_WRITE\_THROUGH}, which returns only once the rename is on disk. All of this uses
the standard library (\code{os}, \code{fcntl}, \code{msvcrt}, \code{ctypes}); no shell command is run.

\begin{table}[t]
\caption{Durability on each system}\label{tab:durable}
\centering\footnotesize
\setlength{\tabcolsep}{4pt}
\begin{tabular}{@{}L{17mm}L{17mm}L{19mm}L{24mm}@{}}
\toprule
\textbf{Step} & \textbf{Linux} & \textbf{macOS} & \textbf{Windows} \\
\midrule
flush a file & \code{fsync} & \code{F\_FULLFSYNC} & \code{FlushFileBuffers} \\ \addlinespace[1.6pt]
flush a folder & \code{fsync} & \code{F\_FULLFSYNC} & not possible; renames written through \\ \addlinespace[1.6pt]
replace a file & \code{rename} & \code{rename} & \code{MoveFileExW}, replace and write-through \\ \addlinespace[1.6pt]
publish without replacing & hard link & hard link & \code{MoveFileExW} without replace \\ \addlinespace[1.6pt]
lock & \code{flock} & \code{flock} & \code{msvcrt.locking} \\
\bottomrule
\end{tabular}
\end{table}

\subsection{The journal}\label{sec:journal}
A record's \code{state} is the journal of the operation in progress (\figref{fig:journal}). It is
written durably before every step that changes something outside the data folder, and every
irreversible step is preceded by a state that tells start-up recovery how to finish or undo it.
The code marks each such point with a named checkpoint; the tests stop the process at every one of
them (\secref{sec:verification}).

\begin{figure}[t]
\centering
% Fig.: the journal states of one ghost record and the transitions between them.
\begin{tikzpicture}[
  font=\scriptsize,
  st/.style={draw=black!65, fill=black!5, rounded corners=5pt, inner xsep=3pt, inner ysep=2.2pt,
             font=\scriptsize\ttfamily, minimum height=5.2mm},
  ok/.style={st, fill=black!78, draw=black!78, text=white},
  gone/.style={draw=black!45, dashed, rounded corners=5pt, inner sep=2.2pt, font=\scriptsize\itshape,
               text=black!65, minimum height=5.2mm},
  arr/.style={-{Latex[length=1.3mm,width=1mm]}, draw=black!65},
  fail/.style={arr, densely dashed},
  lab/.style={font=\scriptsize, text=black!75, align=left,
              execute at begin node={\hyphenpenalty=10000}},
]
\def\step{9.5mm}
\node[st] (up) at (0,0) {uploading};
\node[st, below=\step of up.center, anchor=center] (stored) {stored};
\node[st, below=\step of stored.center, anchor=center] (rem) {removing-original};
\node[ok, below=\step of rem.center, anchor=center] (ready) {ready};
\node[st, below=\step of ready.center, anchor=center] (res) {restoring};
\node[st, below=\step of res.center, anchor=center] (clean) {cleanup};
\node[gone, below=\step of clean.center, anchor=center] (gone) {record removed};
\node[st] (del) at (-29mm, 0 |- ready) {deleting};

\fill[black!70] ([yshift=6mm]up.north) circle (0.9mm) coordinate (start);
\draw[arr] ([yshift=5.1mm]up.north) -- node[lab, right=1mm] {plan and key recorded} (up);
\draw[arr] (up) -- node[lab, right=1mm] {chunks committed and verified} (stored);
\draw[arr] (stored) -- node[lab, right=1mm] {original about to be moved aside} (rem);
\draw[arr] (rem) -- node[lab, right=1mm] {original deleted, or kept} (ready);
\draw[arr] (ready) -- node[lab, left=1mm] {restore} (res);
\draw[arr] (res) -- node[lab, right=1mm] {file published under its name} (clean);
\draw[arr] (clean) -- node[lab, right=1mm] {chunks gone from every branch} (gone);

\draw[fail] (up.west) -| node[lab, pos=0.25, above] {failure} (del.north);
\draw[arr] (ready.west) -- node[lab, above] {Reset} (del.east);
\draw[arr] (del.south) |- node[lab, pos=0.25, right] {chunks gone} (gone.west);
\draw[fail] (res.east) -- ++(4mm,0) |- node[lab, pos=0.25, right] {failure: temporary file removed} (ready.east);
\end{tikzpicture}

\caption{The journal states of a record. Solid arrows are the normal course of an operation;
dashed arrows are taken on failure. An interruption in any state is handled at the next start as
\tabref{tab:recovery} describes.}
\label{fig:journal}
\end{figure}

At every start GhostVault runs this recovery first, as an operation of its own during which the
window accepts no other: it removes leftover temporary files, checks the settings, handles each
record by its state (\tabref{tab:recovery}), and tells the user what it did. The rules
for \state{stored} and \state{removing-original} are written so that a file found at the original's
path is never called the original unless its identity is the uploaded file's: a new file saved under
the same name is not the user's original, and saying so could lead the user to delete it.

\begin{table}[t]
\caption{Start-up recovery by journal state}\label{tab:recovery}
\centering\footnotesize
\begin{tabular}{@{}L{24mm}L{58mm}@{}}
\toprule
\textbf{State} & \textbf{What the next start does} \\
\midrule
\state{uploading} & the record becomes \state{deleting}: every path it lists is removed from the
branches, then the record; the original was never touched \\ \addlinespace[1.6pt]
\state{deleting} & removal from GitHub is finished, then the record is removed \\ \addlinespace[1.6pt]
\state{stored} & the ghost is complete; the original is kept, and the user is told whether the
file at its path is still the original \\ \addlinespace[1.6pt]
\state{removing-original} & if the file under the hidden name exists, it is checked again and
deleted if it is the uploaded file, put back otherwise; the ghost is complete either way \\ \addlinespace[1.6pt]
\state{restoring} & the temporary file is removed; if the final file is in place and matches the
original's SHA-256, it is adopted and clean-up continues, otherwise the ghost is \state{ready}
again \\ \addlinespace[1.6pt]
\state{cleanup} & removal from GitHub is retried; while GitHub cannot be reached it stays pending,
with a Retry button \\
\bottomrule
\end{tabular}
\end{table}

\emph{Unreferenced blobs.} If ghosting is interrupted after some chunks were uploaded but before
they were committed, those blobs exist in the repository's object store without being on any
branch. They are encrypted and listed nowhere, and Git removes unreferenced objects eventually, on
GitHub's schedule. GhostVault cannot delete a blob directly: the Git Data API has no such
operation.

\subsection{Damage}
An unreadable record is shown as damaged and never overwritten, and it pauses backups until the
user resolves it (\secref{sec:incomplete}). An unreadable settings file is moved into
\code{damaged/} and GhostVault starts with empty settings; ghosts are unaffected, because every
record names its own repositories.
